\clearpage
\section*{September 8, 2026: The 99-unit spike across boroughs}
\textit{Agent-drafted entry.}

The post-period spike at 99 units appears in constituent filings in the Bronx,
Brooklyn, Manhattan, and Queens. Aggregating those filings to linked parents
changes the geographic pattern. The Bronx and Brooklyn retain pronounced spikes
at a parent total of 99, while most Manhattan filings at 99 belong to larger
parents. Staten Island has too few observations to support a useful comparison.
The two figures below separate these margins.

At Jacob's request, both comparisons use the existing A/B rental-opportunity
sample in 2019--2022 and January 1, 2025--July 8, 2026. Each borough and period
has its own denominator of observations with at least 50 proposed units,
including the full upper tail. A constituent filing must itself have 50 units
to enter its distribution; parent eligibility uses the sum of its constituent
units. These are unweighted descriptive comparisons. Changes in the composition
of development opportunities and incomplete post-period linkage windows remain
relevant to their interpretation.

Borough comes from the constituent filing records, which agree within every
parent. No observations are excluded because site characteristics are missing.
The saved geographic panel establishes a first reporting baseline: filing keys
are unique, all parent totals reconcile to constituent sums, and the parents
reconstructed here match the canonical parent panel exactly.

The accompanying map locates exact-99 filings and distinguishes parents totaling
99 from larger parents containing a 99-unit filing. It uses historical Housing
Database coordinates and post-period DOB NOW coordinates. One exact-99 filing,
20 Brook Avenue in the Bronx, lacks coordinates and is omitted from the map
while remaining in the borough counts. Coincident locations can overlap. The
map describes locations, not neighborhood bunching rates or causal splitting.

\smallskip
{\small Reproduction: run \texttt{make} in
\texttt{tasks/analyze\_borough\_bunching/code/}. The three-page
\texttt{borough\_bunching.pdf} includes both distributions and the map.
Inputs retain the current cohort definitions, Housing Database 25Q4, and the
existing DOB NOW extract. Borough outlines are the September 8, 2026 NYC Open
Data extract. This entry accompanies working changes on
\texttt{simplify\_empirical\_tasks}, based on revision \texttt{7bbf6d5}.
The current framework is unchanged.}

\clearpage
\begin{center}
\includegraphics[width=\textwidth,height=.92\textheight,keepaspectratio]{../tasks/analyze_borough_bunching/output/borough_filings.pdf}
\end{center}
\clearpage
\begin{center}
\includegraphics[width=\textwidth,height=.92\textheight,keepaspectratio]{../tasks/analyze_borough_bunching/output/borough_parents.pdf}
\end{center}
